\section{Proposed System}

\subsection{Rule-Based RTM Status Calculation}

The RTM module is read-only over existing project data. It does not create or modify
requirements, tasks, tests, bugs, or evidence. Instead, it aggregates linked
artifacts and computes requirement health from task progress, test execution,
bug/defect state, accepted evidence, and risk reasons.

To formalize this logic, let $R$ be a requirement. We define:
\begin{itemize}
    \item $T(R)$: The set of tasks linked to $R$. Let $\text{status}(t) \in \{\text{TODO}, \text{IN\_PROGRESS}, \text{IN\_REVIEW}, \text{DONE}, \text{BLOCKED}\}$ be the status of task $t \in T(R)$.
    \item $TC(R)$: The set of test cases linked to $R$. Let $\text{status}(c) \in \{\text{PASS}, \text{FAIL}, \text{BLOCKED}, \text{NOT\_RUN}\}$ be the latest execution status of test case $c \in TC(R)$.
    \item $B(R)$: The set of bugs linked to $R$ (via related tasks or failed test executions). Let $\text{status}(b) \in \{\text{NEW}, \text{ASSIGNED}, \text{FIXED}, \text{VERIFIED}, \text{CLOSED}\}$ be the status of bug $b \in B(R)$.
    \item $E(R)$: The set of accepted evidences linked to $R$.
    \item $\text{req\_status}(R) \in \{\text{DRAFT}, \text{ACTIVE}, \text{DONE}\}$: The user-defined lifecycle state of $R$.
    \item $\text{evidence\_req}(R) \in \{\text{true}, \text{false}\}$: An indicator of whether evidence is mandatory for $R$.
\end{itemize}

The calculated traceability status of $R$, denoted as $\text{trace\_status}(R) \in \{\text{AT\_RISK}, \text{DONE}, \text{NOT\_STARTED}, \text{IN\_PROGRESS}\}$, is evaluated deterministically using the following hierarchical rules:

\begin{enumerate}
    \item \textbf{At Risk State ($\text{trace\_status}(R) = \text{AT\_RISK}$):} A requirement is flagged as at risk if there are blocked tasks, failed/blocked tests, open bugs, or missing required evidence. Formally:
    \begin{align*}
    \text{trace\_status}(R) = \text{AT\_RISK} \quad \text{if} \quad 
    & \exists t \in T(R): \text{status}(t) = \text{BLOCKED} \\
    & \lor \exists c \in TC(R): \text{status}(c) \in \{\text{FAIL}, \text{BLOCKED}\} \\
    & \lor \exists b \in B(R): \text{status}(b) \notin \{\text{FIXED}, \text{VERIFIED}, \text{CLOSED}\} \\
    & \lor (\text{evidence\_req}(R) \land |E(R)| = 0) \\
    & \lor (T(R) = \emptyset \land \text{req\_status}(R) = \text{ACTIVE})
    \end{align*}
    
    \item \textbf{Done State ($\text{trace\_status}(R) = \text{DONE}$):} A requirement is resolved if it has no blocking risks, its lifecycle status is done, and all linked tasks and tests are successfully completed:
    \begin{align*}
    \text{trace\_status}(R) = \text{DONE} \quad \text{if} \quad 
    & \text{trace\_status}(R) \neq \text{AT\_RISK} \\
    & \land \text{req\_status}(R) = \text{DONE} \\
    & \land (T(R) \neq \emptyset \land \forall t \in T(R): \text{status}(t) = \text{DONE}) \\
    & \land (TC(R) \neq \emptyset \land \forall c \in TC(R): \text{status}(c) = \text{PASS})
    \end{align*}

    \item \textbf{Not Started State ($\text{trace\_status}(R) = \text{NOT\_STARTED}$):} A requirement is untouched if no progress has been made and it has no risks:
    \begin{align*}
    \text{trace\_status}(R) = \text{NOT\_STARTED} \quad \text{if} \quad
    & \text{trace\_status}(R) \neq \text{AT\_RISK} \\
    & \land \forall t \in T(R): \text{status}(t) = \text{TODO} \\
    & \land \forall c \in TC(R): \text{status}(c) = \text{NOT\_RUN} \\
    & \land \text{req\_status}(R) \neq \text{DONE}
    \end{align*}

    \item \textbf{In Progress State ($\text{trace\_status}(R) = \text{IN\_PROGRESS}$):} Otherwise, if there is active development but all criteria for completion have not been met:
    \begin{align*}
    \text{trace\_status}(R) = \text{IN\_PROGRESS} \quad \text{otherwise}
    \end{align*}
\end{enumerate}

\subsection{Evidence-Based Task Review}

DevTrack AI uses a review gate before a task is finally accepted as done. A member
requests review after completing work. The system builds an evidence summary from
available project data, including requirement linkage, assignee information,
checklist status, subtasks, blocked state, linked GitHub issue, uploaded evidence,
and linked code evidence. The project leader then approves or rejects the task. This
process reduces false completion, where a task is marked done without meaningful
proof.

The Code Insight module extends this review gate for code-related tasks. It uses the
shared GitHub Integration Core to handle repository configuration, webhook
verification, and event dispatching. Push, pull request, workflow, and check-run
events can be stored as normalized GitHub evidence and linked to tasks through
explicit task identifiers, GitHub issue references, or SHA chains. Code Insight then
uses local task signals and linked GitHub evidence to produce an explainable score,
show evidence details to the leader, fetch changed-file metadata when needed, store
local structured AI review summaries, capture review snapshots, and expose dashboard
metrics. The leader remains the final authority:

\begin{center}
\texttt{GitHub evidence + deterministic rules + optional AI review + leader decision = trusted task completion}
\end{center}

\subsection{Automated Test Integration via Playwright Service}

Verification of requirements is further automated through online test execution. Rather than relying on manual QA declarations, DevTrack AI integrates an execution service built on Node.js and Playwright. 

When a test run is triggered from the web workspace:
\begin{enumerate}
    \item The Spring Boot backend serializes the test case steps and base URL into a JSON payload and issues an HTTP POST request to the Playwright execution service.
    \item The Playwright service dynamically spins up a headless browser instance, executes the steps sequentially, and records step duration, status, and failures.
    \item Screenshots are captured automatically at the end of execution (or upon failure) and returned as Base64-encoded strings.
    \item The backend parses the response, updates the test execution record, uploads the screenshots to Cloudinary to be cataloged in the Evidence Vault, and recalculates the RTM status.
    \item If the test status is \texttt{FAIL}, the system automatically generates an unassigned Bug Report containing the failed step and step-by-step logs, preventing defects from being overlooked.
\end{enumerate}

\subsection{GitHub Integration and Code Insight Webhook}

To link source-code edits directly to requirements, DevTrack AI utilizes GitHub webhooks. The backend exposes a webhook listener verified via SHA-256 HMAC signatures to guarantee event integrity. 

When a developer pushes code to the repository:
\begin{enumerate}
    \item The GitHub webhook propagates commit and pull request metadata to the DevTrack AI backend.
    \item A commit parser scans commit messages using a regular expression to capture conventional commits linked to tasks:
    \begin{equation}
    \text{Regex Pattern: } \backslash\text{b(feat|fix|docs|refactor|test)}\backslash(\text{TASK-}\backslash\text{d}+)\backslash):\backslash\text{s}+(.+)
    \end{equation}
    For instance, a commit message \texttt{feat(TASK-12): implement login API} is resolved as task key \texttt{TASK-12}.
    \item The commit metadata (SHA, author name, timestamp, and message) is registered as a git evidence record linked to the task.
    \item Commits that do not match the convention are stored as unlinked commits. The project leader can inspect them in the Code Insight dashboard and link them to tasks manually.
\end{enumerate}

\subsection{Factual Grounding and AI-Assisted Reporting}

DevTrack AI provides AI assistance for tasks such as drafting weekly progress reports. To guarantee reports remain accurate and represent actual project velocity, the system uses a factual grounding mechanism. 

The report generation pipeline works as follows:
\begin{enumerate}
    \item The system queries the database to extract structured project facts for the given week: sprint accomplishments, total tasks completed, active blocker warnings, RTM snapshots, latest test run statistics, and red-alert members (members with more than 3 overdue tasks or stale tasks).
    \item These facts are structured into a JSON payload representing the ground truth.
    \item The backend populates a prompt template, merging the structured facts with instructions. A sample prompt structure is shown in Fig.~\ref{fig:prompt}.
    \item The prompt is dispatched to the AI provider (Gemini or OpenAI API). The AI model rewrites the structured metadata into a readable narrative draft.
    \item The project leader reviews the AI-generated draft in the report editor, makes corrections, and saves it. 
\end{enumerate}

This human-in-the-loop design prevents the model from generating hallucinated accomplishments, as any narrative description must match the underlying RTM and task database records.

\begin{figure}
\centering
\begin{minipage}{0.9\textwidth}
\begin{verbatim}
[System Instructions]
You are a project manager writing a weekly report based strictly on facts.
Do not hallucinate. Do not add tasks that are not in the list.

[Factual Project Context]
Project: DevTrack AI | Sprint: Sprint 3
Tasks Completed: 8/10 | Overdue Tasks: 2 (TASK-14, TASK-15)
RTM Health: 85% requirements Done, 15% At Risk (failed test on REQ-2)
Risks Identified: Minh Hieu is flagged RED due to 4 overdue tasks.

[Output Format]
Write a 3-paragraph summary covering: 
1) Completed work, 2) Blockers & risks, 3) Proposed next steps.
\end{verbatim}
\end{minipage}
\caption{Grounding prompt template structure for AI weekly reports}
\label{fig:prompt}
\end{figure}
